\documentclass[main.tex]{subfiles}

\begin{document}

\section{Universal Chern--Weil homomorphism}\label{s:chern}

In this section we recall the construction of a homomorphism 
\begin{equation}\label{eq:cw1}
  \mathsf{cw}_d \colon \C[\op{ch}_1,\ldots,\op{ch}_d]_{d+1} \to \HH^2_{Lie} (\lie{witt}_d) .
\end{equation}
from \cite[section 3]{GWvir2d}.
Here, we treat $\C[\op{ch}_1,\ldots,\op{ch}_d]$ as a graded polynomial ring with generator $\op{ch}_k$ carrying weight
$k$.
In particular, the dimension of its weight $d+1$ subspace is one less than the partition number of the integer $d+1$:
\begin{equation}\label{}
  \dim \C[\op{ch}_1,\ldots,\op{ch}_d]_{d+1} = p(d+1) - 1 
\end{equation}
since we through away the single block partition as we do not include $\op{ch}_{d+1}$.

After recalling the definition of $\mathsf{cw}_d$ we move onto the first main result in this paper which is to show that
$\mathsf{cw}_d$ is injective for all $d$.
The fact that $\op{cw}_1$ is an isomorphism is a classic result in conformal field theory and vertex algebras.
Our proof strategy in this section for general dimension $d$ is a generalization of the argument in \cite{GWvir2d}
for $d=2$.

In the last part of the section we generalize the Chern--Weil homomorphism \ref{eq:cw1} with $\lie{witt}_d$ replaced by
the semi-direct product $\lie{witt}_d \ltimes \lie{gl}_r(\cJ_d)$.

\subsection{The universal Chern--Weil homomorphism}
In this section we will be rather brief and refer to \cite{GWvir2d} for more details and proofs.

\paragraph
There is a universal element 
\[
\lie{c}\in C^\bu_{Lie}\left(\lie{witt}_d ; \br C_\bu^\lambda(\cJ_d) \right)^1
\]
in the Chevalley--Eilenberg cochain complex of $\lie{witt}_d$ with
coefficients in the Connes' reduced cyclic complex computing reduced cyclic homology of $\cJ_d$.
This element is of total degree $+1$ and decomposes as
\begin{equation}\label{eq:c}
  \lie{c} = \lie{c}_1 + \lie{c}_2+ \cdots 
\end{equation}
where
\begin{equation}\label{}
  \lie{c}_{k+1} \colon \wedge^{k+1} \lie{witt}_d \to \br C_\bu^\lambda(\cJ_d)
\end{equation}
is defined by
\begin{align}
  \lie{c}_{k+1}(T_0\wedge\cdots\wedge T_{k}) &=\frac{(-1)^k}{2^k} \sum_{\permute \in S_k}(\pm)^{'} \cdot \op{Tr}\left(J T_{0}, J T_{\permute(1)}, \ldots, J T_{\permute(k)}\right) \\
                                             & =\frac{(-1)^k}{2^k\cdot(k+1)} \sum_{\permute' \in S_{k+1}}(\pm)^{''}
                                             \cdot \op{Tr}\left(J T_{\permute'(0)}, J T_{\permute'(1)}, \ldots, J T_{\permute'(k)}\right) 
\end{align}
where
$$
(\pm)'=\mathrm{sgn}(\permute)\cdot\varepsilon(\permute;T_1,\ldots,T_k) ,
$$
$$
(\pm)''=\mathrm{sgn}(\permute')\cdot\varepsilon(\permute';T_0,T_1,\ldots,T_k) ,
$$
and $J \colon \lie{witt}_d \to \lie{gl}_d(\cJ_d)$ is defined in equation \eqref{eq:jacobian}. 

For $A$ a commutative dg algebra, Connes' complex $\overline{C}_\bu^\lambda(A)$ admits, itself, the natural structure of a
commutative dg algebra \cite[proposition 3.7]{loday1984cyclic}. 
It is defined by the formula
\begin{equation}\label{eqn:starproduct}
x*y \define x\times(By),
\end{equation}
where
\[
B(a_0,a_1,\dots,a_d)=\sum^n_{i=0}\pm (1,a_i,\dots,a_n,a_0,\dots,a_{i-1}),
\]
and 
\[
\pm=(-1)^{i\cdot n}\cdot\varepsilon(\permute_i;a_0,a_1,\dots,a_d),\quad \permute_i(a_0,a_1,\dots,a_d)=(a_i,\dots,a_n,a_0,\dots,a_{i-1})
\]
The operation $\times$ is the so-called shuffle product, given by 
\[
(a,a_1,\dots,a_p)\times (a',a_{p+1},\dots,a_{p+q})=\sum_{\text{Shuffle}}(\pm)\cdot\mathrm{sgn}(\permute)\cdot (aa',a_{\permute^{-1}(1)},\dots,a_{\permute^{-1}(p+q)}),
\]
\[
(\pm)=(-1)^{|a'|\sum^p\limits_{i=1}|a_i|}\cdot\varepsilon(\permute_i;a_1,\dots,a_{p+q})
\]
where the sum is over all permutations $\permute$ of $\{1,\dots,p+q\}$ such that $\permute(1)<\cdots<\permute(p)$ and $\permute(p+1)<\cdots<\permute(p+q)$.

Define the cyclic cocycle~$\rho \in C^\bu_\lambda(\cJ_d)$ of total degree one by the formula
\begin{equation}\label{eq:rho}
  \rho(f_0,\ldots,f_d) = \op{Res}(f_0 \del f_1 \cdots \del f_d) .
\end{equation}
where $\op{Res}(- \d^d z)$ is the unique $GL_d$-invariant degree $-d$ cochain map $\cJ_d \to \C$ with the property that $\op{Res}
(\d^d z) =1$.

\paragraph
We are now ready to define the universal Chern--Weil homomorphism
\begin{equation}\label{}
  \mathsf{cw}_d \colon \C[y_1,\ldots,y_d]_{d+1} \to \HH^2_{Lie}(\lie{witt}_d) ,
\end{equation}
where the $y_1,\ldots,y_d$ are polynomial variables and $y_i$ carries homogenous degree $i$.
The subscript in the domain means that we look at total degree $d+1$ elements.
Let $i_1,\ldots,i_d$ be integers such that $i_1+2i_2+\cdots+d i_d=d+1$.
Define
\begin{equation}\label{eq:chernwitt}
  \mathsf{cw}_d (y_1^{i_1} y_2^{i_2} \cdots y_d^{i_d}) \define \rho \left(\lie{c}_1^{*i_1} * \lie{c}_2^{*i_2} * \cdots
    * \lie{c}_d^{i_d}\right) .
\end{equation}
Implicit in this definition is the claim that the this element is in fact a cocycle and we are taking its cohomology
class.

\paragraph[p:cwdr]
The above construction naturally generalizes to the algebra $\lie{witt}_d\ltimes \lie{gl}_r(\cJ_d)$ in the following way.
Suppose that 
\[
  \theta \in \op{S}^{k+1}(\lie{gl}_r^*)^{\lie{gl}_r}
\]
is an invariant polynomial on $r \times r$ matrices of
homogeneous polynomial degree $k+1$.
Define
\begin{equation}\label{}
  \gamma_{\theta} \colon \wedge^{k+1} \lie{gl}_r(\cJ_d) \to \br C_\bu^\lambda(\cJ_d)[1]
\end{equation}
by the formula
\begin{align}
  \gamma_{\theta}(M_0\otimes f_0\wedge\cdots\wedge M_k\otimes f_k) &=\frac{(-1)^k}{2^k} \sum_{\permute \in S_k}(\pm)^{'}
  \cdot \theta(M_0,M_{\permute(1)},\dots,M_{\permute(k)})\cdot \left(f_0, f_{\permute(1)}, \ldots, f_{\permute(k)}\right) .
\end{align}
It follows immediately that $\gamma_\theta$ is a cochain map, $\lie{gl}_r$-equivariant, and also $\lie{witt}
_d$-equivariant.

\paragraph
Let $\theta$ be a degree $k$ invariant polynomial on $r \times r$ matrices as above and suppose $i_1,\ldots,i_{d-k}$ are integers such that 
\begin{equation}\label{}
  i_1 + 2 i_2 + \cdots + (d-k) i_{d-k} = d-k . 
\end{equation}
Define
\begin{equation}\label{}
  \mathsf{cw}_{d,r} (y_1^{i_1} \cdots y_{d-k}^{i_{d-k}} \theta) \define \mathsf{cw}\left( \lie{c}_1^{*i_1} * \cdots
  *\lie{c}_{d-k}^{*i_{d-k}} * \gamma_\theta \right) ,
\end{equation}
which is an element in $\HH^2_{Lie}(\lie{witt}_d \ltimes \lie{gl}_r(\cJ_d))$.

Consider now the ring 
\begin{equation}\label{}
  \C[y_1,\ldots,y_d] \otimes \op{S}^\bu(\lie{gl}_r^*)^{\lie{gl}_r} \cong \C[y_1,\ldots,y_d, \lie{t}_1,\ldots, \lie{t}_r] 
\end{equation}
which we assign a non-cohomological weight grading by declaring that $y_i$ is degree $i$ and $\lie{t}_{j} \in \op{S}^j(\lie{gl}^*_r)$
defined by 
\begin{equation}\label{}
  \lie{t}_j (A) = \op{tr} (A^j)
\end{equation}
is of degree $j$.
Our construction defines a homomorphism
\begin{equation}\label{}
  \mathsf{cw}_{d,r} \colon \C[y_1,\ldots,y_d, \lie{t}_1,\ldots,\lie{t}_r]_{d+1} \to \HH^2_{Lie}(\lie{witt}_d \ltimes
  \lie{gl}_r(\cJ_d)) 
\end{equation}
Explicitly, on the trace generators, this is defined as
\begin{equation}\label{eq:chernwittMixed}
  \mathsf{cw}_{d,r} (y_1^{i_1} y_2^{i_2} \cdots y_d^{i_d} \cdot \lie{t}_1^{k_1} \cdots \lie{t}_r^{k_r} ) \define \rho \left(\lie{c}_1^{*i_1} * \lie{c}_2^{*i_2} * \cdots
  * \lie{c}_d^{*i_d}* \gamma_{\lie{t}_1}^{* k_1} * \cdots \gamma_{\lie{t}_r}^{* k_r} \right) 
\end{equation}
where $i_1,\ldots,i_d,k_1,\ldots,k_r$ are integers satisfying the relation 
\begin{equation}\label{}
  i_1+2i_2+\cdots+d i_d+k_1 + 2 k_2 + \cdots + r k_r =d+1 .
\end{equation}

\subsection{Injectivity}

We will show (a) the cohomology class of any cocycle of the form \eqref{eq:chernwitt} is nontrivial
and (b) such classes are linearly independent.
In other words, we prove the map $\mathsf{cw}_d$
is injective.
Of course, this statement is obvious and well-known in the case $d=1$.
In \cite{GWvir2d} have proven injectivity when $d=2$.
We treat the general case here.

Our notation in this section is 
\begin{equation}\label{}
  \op{ch}_1^{i_1} \cdots \op{ch}_d^{i_d} \define \mathsf{cw}_d (y_1^{i_1} \cdots y_d^{i_d}) 
\end{equation}
where $\sum^{d}\limits_{p=2} p i_p = d+1$.

\paragraph
We begin with the following lemma which expresses the top trace class as a linear combination of elements in the image
of $\mathsf{cw}_d$.

\begin{lemma}\label{lem:chd+1}
In $\mathbb{H}^2_{\mathrm{Lie}}\left(\lie{witt}_{d}\right)$, the class 
\[
  \mathrm{ch}_{d+1} \define \rho(\lie{c}_{d+1}) 
\]
can be expressed as linear combination of elements of the form 
\[
  (\mathrm{ch}_1)^{i_1}(\mathrm{ch}_2)^{i_2}\cdots (\mathrm{ch}_d)^{i_d}
\]
where $\sum^{d}\limits_{p=1}p i_p=d+1$.
In other words, $\op{ch}_{d+1}$ is in the image of $\mathsf{cw}_d$.
\end{lemma}
\begin{proof}
    The proof is by the Procesi-Razmyslov fundamental trace identity, equivalently the multilinear polarized
    Cayley-Hamilton trace identity (see for example \cite[Theorem 4.3(b)]{PROCESI1976306}). We apply the following trace identity to the matrix ring $\lie{gl}_d(\cJ_d^{\otimes d+1})$
    \[
    \sum_{\tau\in S_{d+1}}\mathrm{sgn}(\tau)\prod_{C=(i_1...i+r)\in \mathrm{Cycles}(\tau)}\mathrm{Tr}\big(M_{i_1}\cdots M_{i_r}\big)=0,\quad M_0,...,M_d\in \lie{gl}_d(\cJ_d^{\otimes d+1}).
    \]
    The explicit formula is 
    \[
      \op{ch}_{d+1} =
\frac{(-1)^{d+1}}{d!}
\op{exp}\left(
  \sum_{k=1}^d (-1)^{k-1}(k-1)!\op{ch}_k t^k
\right) |_{t^{d+1}} 
\]
where the subscript means we take the coefficient of $t^{d+1}$.
\end{proof}

\paragraph
In what follows it will be more convenient to label our degree $d+1$ classes by partitions of the integer $d+1$.
For example 
\[
  \op{ch}_{(1,\ldots,1)} = \op{ch}_1^{d+1} 
\]
while
\[
  \op{ch}_{(2,1,\ldots,1)} = \op{ch}_1^{d-1} \op{ch}_2 . 
\]
More generally, if $\lambda = (1^{i_1}, \ldots, d^{i_d})$ is a partition of $d+1$ then we define 
\begin{equation}\label{}
  \op{ch}_\lambda \define \op{ch}_1^{i_1} \cdots \op{ch}_d^{i_d}
\end{equation}
We also set $\op{ch}_{(d+1)} = \op{ch}_{d+1}$.

\paragraph{Injectivity in dimension $d=2$}
Before diving into the main proof of injectivity, we survey the $d=2$ case. 
When $d=2$ we have two classes $\op{ch}_1^3$ and $\op{ch}_1 \op{ch}_2$ which span $\HH^2_{Lie}(\lie{witt}_d)$.
Our strategy is to find, for example, a chain $\til{\mathbb{v}}\in \bigwedge^{\bullet}\lie{witt}_{2}$ of total degree $-2$
which detects the nontriviality of $\op{ch}_1^3$ at the
cochain level:
$$
(\dbar + \bd^{(2)})\til{\mathbb{v}}=0,\quad (\mathrm{ch}_1^3)(\til{\mathbb{v}})\neq 0.
$$
Here, $\dbar$ is the linear part of the differential and $\bd^{(2)}$ is induced from the Lie bracket.
In other words, $\til{\mathbb{v}}$ is a cycle in the Chevalley--Eilenberg complex for $\lie{witt}_2$, and it pairs
nontrivially with the cocycle $\op{ch}_1^3$.
This shows that the image of $\mathsf{cw}_2$ is at least one-dimensional.
To show that it is two-dimensional, it suffices to find another chain which pairs nontrivially with a linearly
independent cohomology class.
One obvious choice is the class $\op{ch}_1 \op{ch}_2$.
Instead, we will find it easier to find a cycle which is dual to the class $\op{ch}_3$, which, in this dimension, is a linear combination 
\begin{equation}\label{}
  \op{ch}_3 = -\frac{1}{12} \op{ch}_1^3 + \frac12  \op{ch}_1 \op{ch}_2 .
\end{equation}


In dimension two, it is easy to write down the cycles which implement this strategy.
One starts with the two chains:
\begin{align*}
  \mathbb{v}_{(1,1,1)}(2) & \define  2 P_2\partial_{\1}\wedge (z^{\1})^3\partial_{\1}\wedge (z^{\2})^2\partial_{\2} \\
  \mathbb{v}_{(3)}(2) & \define P_2\partial_{\1}\wedge z^{\1}z^{\1}\mathsf{Eu}\wedge z^{\2} \mathsf{Eu}
  -P_2\partial_{\1} \wedge z^{\1}z^{\2} \mathsf{Eu}\wedge z^{\1} \mathsf{Eu},
\end{align*}
where $\mathsf{Eu} = z^\1 \del_{\1} + z^{\2} \del_{\2}$ is the total Euler vector field.
We will explain momentarily, see definition \ref{dfn:partitionv}, how one associates a chain of this type to a general
partition.
The key point is that these chains pair nontrivially with $\op{ch}_1^3$ and $\op{ch}_3$ respectively, see table
\ref{tbl:pairing2} (and ignore the column $\op{ch}_{(2,1)}(2)$ which we have included only for transparency).

\begin{table}[!h]
        \centering
        
\begin{tabular}{|p{0.10\textwidth}|p{0.12\textwidth}|p{0.10\textwidth}|p{0.10\textwidth}|}
\hline 
  & $\displaystyle \mathrm{ch}_{(1,1,1)}(2)$ & $\displaystyle \mathrm{ch}_{(2,1)}(2)$ & $\displaystyle \mathrm{ch}_{(3)}(2)$ \\
\hline 
 $\displaystyle \mathbb{v}_{(1,1,1)}( 2)$ & $\displaystyle \neq 0$ & $\displaystyle \neq 0$ & $\displaystyle 0$ \\
\hline 
 $\displaystyle \mathbb{v}_{(3)}( 2)$ & $\displaystyle \neq 0$ & $\displaystyle \neq 0$ & $\displaystyle \neq 0$ \\
 \hline
\end{tabular}
\caption{Pairing matrix, $d=2$}
\label{tbl:pairing2}
        \end{table}

On the other hand, the chains above are not cycles in the complex $C_\bu^{Lie}(\lie{witt}_d)$ computing the Lie
algebra hyperhomology of $\lie{witt}_d$.
Fortunately, we can run a version of descent to introduce deformations of these chains which are cocycles.
Explicitly, in dimension $d=2$, one choice of descent data leads to:
\begin{align*}
\til{\mathbb v}_{(1,1,1)}(2)
&=
2P_2\partial_{\1}\wedge (z^{\1})^3\partial_{\1}\wedge (z^{\2})^2\partial_{\2}
\\
&\quad
+
\left(
7z^{\1}x^{\2}
+
  (z^{\1})^2x^{\1}x^{\2}
\right)
\partial_{\1}\wedge (z^{\2})^2\partial_{\2}
\\
&\quad
+
\left(
x^{\1}
+
z^{\2}x^{\1}x^{\2}
\right)
\partial_{\1}\wedge (z^{\1})^3\partial_{\1}
\\
&\quad
-
\left(
3(z^{\1})^2(x^{\1})^2
-
6z^{\1}x^{\1}
+
7
\right)
z^{\1}x^{\2}\,
\partial_{\1}\wedge z^{\2}\partial_{\2}.
\end{align*}
and
\begin{align*}
  \til{\mathbb v}_{(3)}(2)
&=
P_2\partial_{\1}\wedge (z^{\1})^2\mathsf{Eu}\wedge z^{\2}\mathsf{Eu}
-
P_2\partial_{\1}\wedge z^{\1}z^{\2}\mathsf{Eu}\wedge z^{\1}\mathsf{Eu}
\\
&\quad
+
\left(4z^{\1}x^{\2}\partial_{\1}-2z^{\1}x^{\1}\partial_{\2}\right)\wedge z^{\2}\mathsf{Eu}
\\
&\quad
+
2x^{\1}\partial_{\1}\wedge (z^{\1})^2\mathsf{Eu}
\\
&\quad
+
\left(
-z^{\2}x^{\2}\partial_{\1}
+
2z^{\1}x^{\1}\partial_{\1}
+
z^{\2}x^{\1}\partial_{\2}
\right)
\wedge z^{\1}\mathsf{Eu}
\\
&\quad
+
\left(3x^{\2}\partial_{\1}-x^{\1}\partial_{\2}\right)\wedge z^{\1}z^{\2}\mathsf{Eu}
\\
&\quad
-
\left(
-2z^{\1}x^{\1}+11
\right)x^{\2}\,
\partial_{\1}\wedge z^{\1}\mathsf{Eu}
-
2x^{\2}\,
\partial_{\2}\wedge z^{\1}z^{\2}\mathsf{Eu}.
\end{align*}


The last remark is that these correction terms automatically pair trivially, by type reasons, with $\op{ch}_{(1,1,1)}$ and $\op{ch}_{(3)}
$.
Thus, the contents of table \ref{tbl:pairing2} are unchanged if we replace the $\mathbb{v}$'s with their cocycle
corrections $\til{\mathbb{v}}$'s.
        This completes the proof that $\mathsf{cw}_2$ is injective. 

\paragraph{Injectivity for $d\geq 2$}
Recall the Lie algebra of formal vector fields $\lie{w}_d$ on the $d$-disk is a dg Lie subalgebra of $\lie{witt}_d$.
Let $\lie{w}^{poly}_d \subset \lie{w}_d$ denote the Lie subalgebra of formal vector fields whose coefficients are
polynomial in $d$-variables.

We work with Chevalley--Eilenberg chains rather than cochains.
We let $\dbar + \bd^{(2)}$ denote the total differential acting on the Chevalley--Eilenberg complex
$C_\bu(\lie{witt}^{poly}_d)$ computing the Lie algebra homology of $\lie{witt}^{poly}_d$.
(So, $\dbar$ is the linear part of the differential and $\bd^{(2)}$ is induced from the bracket.)
In the following definition, we view $\wedge^\bu \lie{w}_d^{poly}$ as subspace of the total degree $(-d)$ part of
$C_\bu(\lie{witt}^{poly}_d)$.

\begin{definition}
  Let $\mathbf{Z}_1(d)\subset \wedge^d \lie{w}^{poly}_d$ be the subspace consisting of
    $$
    \mathbb{w}\in \wedge^d \lie{w}^{\mathrm{pol}}_d,\quad \mathbf{d}^{(2)}\mathbb{w}=0,\ \text{and}\  \mathbf{d}^{(2)}
    (P_d\partial_{\1}\wedge\mathbb{w})=\dbar(-).
    $$
    In other words, $d$-cochains $\mathbb{w}$ which are closed for the Chevalley--Eilenberg differential of formal vector fields and
    such that in the full Chevalley--Eilenberg complex for $\lie{witt}_d$ the quadratic part of the
    Chevalley--Eilenberg applied to $P_d \del_{\1} \wedge \mathbb{w}$ is $\dbar$-exact.
\end{definition}

%Recall the following classical fact in \cite[pp84,Lemma 2.]{FuksBook} (we reformulate in terms of homology by using the
%fact that the continuous linear dual of $H_{q}(\lie{w}_d^{poly})$ is isomorphic to $H^q(\lie{w}_d)$.
%\begin{lemma}
%  $H_{q}(\lie{w}_d^{poly})=0$ for $0<q\leq2n$.
%\end{lemma}
%This vanishing result easily leads to the following.
%Let $F_r C_\bu^{Lie}(\lie{witt}_d)$ be the increasing filtration which consists of chains which carry total graded
%polynomial degree $\leq r$.

\begin{lemma}\label{lem:descent}
  For $\mathbb{w}\in \mathbf{Z}_1(d)$, the class of the cochain 
  \[
    \mathbb{v}^{(d-1)} \define P_d\del_{\1} \wedge\mathbb{w}
  \]
  in $\op{gr}_{d+1} C^{Lie}_{\bu}(\lie{witt}_d)$ lifts
  to a cocycle in $F_{d+1} C^{Lie}_{\bu}(\lie{witt}_d)$.
\end{lemma}
\begin{proof}
    We refer to figure \ref{fig:Zig-Zag} to elucidate the following zig-zag argument.
    Starting with $\mathbb{v}^{(d-1)}$, we construct a total degree $(-2)$ cocycle in $C^{Lie}_\bu (\lie{witt}_d)$ of the form 
    \begin{equation}\label{eq:sumcocycle}
      \til{\mathbb{v}} = \mathbb{v}^{(d-1)} + \cdots + \mathbb{v}^{(0)} 
    \end{equation}
where $\mathbb{v}^{(j)}$ has internal (Joaunolou) degree $j$.
    By assumption, there is $\mathbb{v}^{(d-2)}$ such that 
    \[
      \dbar \mathbb{v}^{(d-2)} + \bd^{(2)} \mathbb{v}^{(d-1)} = 0 .
    \]
    From this equation, we see that $\dbar\bd^{(2)} \mathbb{v}^{(d-2)} = 0$.
    Since the cohomology of the Joaunolou model is concentrated in degrees up to $0$ and $d-1$ we see that there exists
    $\mathbb{v}^{(d-3)}$ such that $\dbar \mathbb{v}^{(d-3)} + \bd^{(2)} \mathbb{v}^{(d-2)} = 0$.
    Continuing on in this way, we produce $\mathbb{v}^{(d-1)}, \mathbb{v}^{(d-2)}, \ldots, \mathbb{v}^{(1)}, \mathbb{w}$
    together which
    satisfy the descent equations 
    \begin{align*}
      \dbar \mathbb{v}^{(d-1)} & = 0 \\
      \dbar \mathbb{v}^{(d-2)} + \bd^{(2)} \mathbb{v}^{(d-1)} & = 0 \\
                                                             & \vdots \\
      \dbar \mathbb{w} + \bd^{(2)} \mathbb{v}^{(1)} & = 0 .
    \end{align*}
    In particular:
    \begin{equation}\label{}
      \left(\dbar + \bd^{(2)}\right)\left(\mathbb{v}^{(d-1)} + \cdots + \mathbb{v}^{(1)} + \mathbb{w}\right) = \bd^{(2)} 
      \mathbb{w} .
    \end{equation}
    In other words, the obstruction for the sum inside parentheses to be a cocycle is $\bd^{(2)} \mathbb{w}$.
    The issue is that we cannot assume that $\mathbb{w}$ is an element of $\wedge^2 \lie{w}_d^{poly}$, it is an element
    of $\wedge^2 (\lie{witt}_d^{poly})^0$.
    But, from the last descent equation we see that $\bd^{(2)} \mathbb{w}$ is $\dbar$-closed hence 
    \begin{equation}\label{}
      \bd^{(2)} \mathbb{w} \in \lie{w}^{poly}_d . 
    \end{equation}
    Since this element is automatically $\bd^{(2)}$-closed (where, we now mean the Chevalley--Eilenberg differential for
    $\lie{w}_d^{poly}$), it follows from $H_1^{Lie}(\lie{w}^{poly}_d) = 0$
    that there exists $\mathbb{u} \in \wedge^2 \lie{w}_d$ with the property that $\bd^{(2)} \mathbb{u} =
    \bd^{(2)} \mathbb{w}$.
    Finally, define:
    \begin{equation}\label{}
      \mathbb{v}^{(0)} \define \mathbb{w} - \mathbb{u} .
    \end{equation}
    The desired cocycle is then \eqref{eq:sumcocycle}.
    \begin{figure}[htbp]
        \centering

\tikzset{every picture/.style={line width=0.75pt}} %set default line width to 0.75pt        

\begin{tikzpicture}[x=0.75pt,y=0.75pt,yscale=-1,xscale=1]
%uncomment if require: \path (0,362); %set diagram left start at 0, and has height of 362

%Shape: Axis 2D [id:dp4320338525459244] 
\draw  (475.2,323.56) -- (474.49,35.63)(125.93,295.62) -- (513.92,294.67) (469.51,42.64) -- (474.49,35.63) -- (479.51,42.62) (132.94,300.6) -- (125.93,295.62) -- (132.92,290.6) (470,243.78) -- (480,243.76)(469.88,192.78) -- (479.88,192.76)(469.75,141.78) -- (479.75,141.76)(469.63,90.78) -- (479.63,90.76)(424.14,299.89) -- (424.11,289.89)(373.14,300.02) -- (373.11,290.02)(322.14,300.14) -- (322.11,290.14)(271.14,300.27) -- (271.11,290.27)(220.14,300.39) -- (220.11,290.39)(169.14,300.52) -- (169.11,290.52) ;
\draw   ;
%Straight Lines [id:da2900421336009549] 
\draw [color={rgb, 255:red, 189; green, 16; blue, 224 }  ,draw opacity=1 ][fill={rgb, 255:red, 74; green, 144; blue, 226 }  ,fill opacity=1 ]   (221,93) ;
\draw [shift={(221,93)}, rotate = 0] [color={rgb, 255:red, 189; green, 16; blue, 224 }  ,draw opacity=1 ][fill={rgb, 255:red, 189; green, 16; blue, 224 }  ,fill opacity=1 ][line width=0.75]      (0, 0) circle [x radius= 3.35, y radius= 3.35]   ;
%Straight Lines [id:da5500088828853273] 
\draw [color={rgb, 255:red, 0; green, 0; blue, 0 }  ,draw opacity=1 ][fill={rgb, 255:red, 74; green, 144; blue, 226 }  ,fill opacity=1 ]   (208,92.91) -- (169,92.91) ;
\draw [shift={(169,92.91)}, rotate = 180] [color={rgb, 255:red, 0; green, 0; blue, 0 }  ,draw opacity=1 ][fill={rgb, 255:red, 0; green, 0; blue, 0 }  ,fill opacity=1 ][line width=0.75]      (0, 0) circle [x radius= 3.35, y radius= 3.35]   ;
\draw [shift={(210,92.91)}, rotate = 180] [fill={rgb, 255:red, 0; green, 0; blue, 0 }  ,fill opacity=1 ][line width=0.08]  [draw opacity=0] (12,-3) -- (0,0) -- (12,3) -- cycle    ;
%Straight Lines [id:da1503395154909084] 
\draw [color={rgb, 255:red, 0; green, 0; blue, 0 }  ,draw opacity=1 ][fill={rgb, 255:red, 74; green, 144; blue, 226 }  ,fill opacity=1 ]   (221.56,102) -- (221.27,141) ;
\draw [shift={(221.27,141)}, rotate = 90.44] [color={rgb, 255:red, 0; green, 0; blue, 0 }  ,draw opacity=1 ][fill={rgb, 255:red, 0; green, 0; blue, 0 }  ,fill opacity=1 ][line width=0.75]      (0, 0) circle [x radius= 3.35, y radius= 3.35]   ;
\draw [shift={(221.58,100)}, rotate = 90.44] [fill={rgb, 255:red, 0; green, 0; blue, 0 }  ,fill opacity=1 ][line width=0.08]  [draw opacity=0] (12,-3) -- (0,0) -- (12,3) -- cycle    ;
%Straight Lines [id:da014870988023802267] 
\draw [color={rgb, 255:red, 189; green, 16; blue, 224 }  ,draw opacity=1 ][fill={rgb, 255:red, 74; green, 144; blue, 226 }  ,fill opacity=1 ]   (273,141) ;
\draw [shift={(273,141)}, rotate = 0] [color={rgb, 255:red, 189; green, 16; blue, 224 }  ,draw opacity=1 ][fill={rgb, 255:red, 189; green, 16; blue, 224 }  ,fill opacity=1 ][line width=0.75]      (0, 0) circle [x radius= 3.35, y radius= 3.35]   ;
%Straight Lines [id:da2043361360051017] 
\draw [color={rgb, 255:red, 0; green, 0; blue, 0 }  ,draw opacity=1 ][fill={rgb, 255:red, 74; green, 144; blue, 226 }  ,fill opacity=1 ]   (260,140.91) -- (221,140.91) ;
\draw [shift={(221,140.91)}, rotate = 180] [color={rgb, 255:red, 0; green, 0; blue, 0 }  ,draw opacity=1 ][fill={rgb, 255:red, 0; green, 0; blue, 0 }  ,fill opacity=1 ][line width=0.75]      (0, 0) circle [x radius= 3.35, y radius= 3.35]   ;
\draw [shift={(262,140.91)}, rotate = 180] [fill={rgb, 255:red, 0; green, 0; blue, 0 }  ,fill opacity=1 ][line width=0.08]  [draw opacity=0] (12,-3) -- (0,0) -- (12,3) -- cycle    ;
%Straight Lines [id:da7612997244093023] 
\draw [color={rgb, 255:red, 0; green, 0; blue, 0 }  ,draw opacity=1 ][fill={rgb, 255:red, 74; green, 144; blue, 226 }  ,fill opacity=1 ]   (273.56,150) -- (273.27,189) ;
\draw [shift={(273.27,189)}, rotate = 90.44] [color={rgb, 255:red, 0; green, 0; blue, 0 }  ,draw opacity=1 ][fill={rgb, 255:red, 0; green, 0; blue, 0 }  ,fill opacity=1 ][line width=0.75]      (0, 0) circle [x radius= 3.35, y radius= 3.35]   ;
\draw [shift={(273.58,148)}, rotate = 90.44] [fill={rgb, 255:red, 0; green, 0; blue, 0 }  ,fill opacity=1 ][line width=0.08]  [draw opacity=0] (12,-3) -- (0,0) -- (12,3) -- cycle    ;
%Straight Lines [id:da18437113410659078] 
\draw [color={rgb, 255:red, 0; green, 0; blue, 0 }  ,draw opacity=1 ][fill={rgb, 255:red, 74; green, 144; blue, 226 }  ,fill opacity=1 ]   (313,188.91) -- (274,188.91) ;
\draw [shift={(274,188.91)}, rotate = 180] [color={rgb, 255:red, 0; green, 0; blue, 0 }  ,draw opacity=1 ][fill={rgb, 255:red, 0; green, 0; blue, 0 }  ,fill opacity=1 ][line width=0.75]      (0, 0) circle [x radius= 3.35, y radius= 3.35]   ;
\draw [shift={(315,188.91)}, rotate = 180] [fill={rgb, 255:red, 0; green, 0; blue, 0 }  ,fill opacity=1 ][line width=0.08]  [draw opacity=0] (12,-3) -- (0,0) -- (12,3) -- cycle    ;
%Straight Lines [id:da18159335931665033] 
\draw [color={rgb, 255:red, 189; green, 16; blue, 224 }  ,draw opacity=1 ][fill={rgb, 255:red, 74; green, 144; blue, 226 }  ,fill opacity=1 ]   (321,192) ;
\draw [shift={(321,192)}, rotate = 0] [color={rgb, 255:red, 189; green, 16; blue, 224 }  ,draw opacity=1 ][fill={rgb, 255:red, 189; green, 16; blue, 224 }  ,fill opacity=1 ][line width=0.75]      (0, 0) circle [x radius= 3.35, y radius= 3.35]   ;
%Straight Lines [id:da27425030392278593] 
\draw [color={rgb, 255:red, 0; green, 0; blue, 0 }  ,draw opacity=1 ][fill={rgb, 255:red, 74; green, 144; blue, 226 }  ,fill opacity=1 ]   (320.56,203) -- (320.27,242) ;
\draw [shift={(320.27,242)}, rotate = 90.44] [color={rgb, 255:red, 0; green, 0; blue, 0 }  ,draw opacity=1 ][fill={rgb, 255:red, 0; green, 0; blue, 0 }  ,fill opacity=1 ][line width=0.75]      (0, 0) circle [x radius= 3.35, y radius= 3.35]   ;
\draw [shift={(320.58,201)}, rotate = 90.44] [fill={rgb, 255:red, 0; green, 0; blue, 0 }  ,fill opacity=1 ][line width=0.08]  [draw opacity=0] (12,-3) -- (0,0) -- (12,3) -- cycle    ;
%Straight Lines [id:da5538174997058338] 
\draw [color={rgb, 255:red, 189; green, 16; blue, 224 }  ,draw opacity=1 ][fill={rgb, 255:red, 74; green, 144; blue, 226 }  ,fill opacity=1 ]   (372,242) ;
\draw [shift={(372,242)}, rotate = 0] [color={rgb, 255:red, 189; green, 16; blue, 224 }  ,draw opacity=1 ][fill={rgb, 255:red, 189; green, 16; blue, 224 }  ,fill opacity=1 ][line width=0.75]      (0, 0) circle [x radius= 3.35, y radius= 3.35]   ;
%Straight Lines [id:da9487501648592407] 
\draw [color={rgb, 255:red, 0; green, 0; blue, 0 }  ,draw opacity=1 ][fill={rgb, 255:red, 74; green, 144; blue, 226 }  ,fill opacity=1 ]   (359,241.91) -- (320,241.91) ;
\draw [shift={(320,241.91)}, rotate = 180] [color={rgb, 255:red, 0; green, 0; blue, 0 }  ,draw opacity=1 ][fill={rgb, 255:red, 0; green, 0; blue, 0 }  ,fill opacity=1 ][line width=0.75]      (0, 0) circle [x radius= 3.35, y radius= 3.35]   ;
\draw [shift={(361,241.91)}, rotate = 180] [fill={rgb, 255:red, 0; green, 0; blue, 0 }  ,fill opacity=1 ][line width=0.08]  [draw opacity=0] (12,-3) -- (0,0) -- (12,3) -- cycle    ;
%Straight Lines [id:da43947256340248986] 
\draw [color={rgb, 255:red, 0; green, 0; blue, 0 }  ,draw opacity=1 ][fill={rgb, 255:red, 74; green, 144; blue, 226 }  ,fill opacity=1 ]   (372.56,251) -- (372.27,290) ;
\draw [shift={(372.27,290)}, rotate = 90.44] [color={rgb, 255:red, 0; green, 0; blue, 0 }  ,draw opacity=1 ][fill={rgb, 255:red, 0; green, 0; blue, 0 }  ,fill opacity=1 ][line width=0.75]      (0, 0) circle [x radius= 3.35, y radius= 3.35]   ;
\draw [shift={(372.58,249)}, rotate = 90.44] [fill={rgb, 255:red, 0; green, 0; blue, 0 }  ,fill opacity=1 ][line width=0.08]  [draw opacity=0] (12,-3) -- (0,0) -- (12,3) -- cycle    ;
%Straight Lines [id:da0649212091209419] 
\draw [color={rgb, 255:red, 189; green, 16; blue, 224 }  ,draw opacity=1 ][fill={rgb, 255:red, 74; green, 144; blue, 226 }  ,fill opacity=1 ]   (424,290) ;
\draw [shift={(424,290)}, rotate = 0] [color={rgb, 255:red, 189; green, 16; blue, 224 }  ,draw opacity=1 ][fill={rgb, 255:red, 189; green, 16; blue, 224 }  ,fill opacity=1 ][line width=0.75]      (0, 0) circle [x radius= 3.35, y radius= 3.35]   ;
%Straight Lines [id:da2438265525064589] 
\draw [color={rgb, 255:red, 0; green, 0; blue, 0 }  ,draw opacity=1 ][fill={rgb, 255:red, 74; green, 144; blue, 226 }  ,fill opacity=1 ]   (412,289.91) -- (373,289.91) ;
\draw [shift={(373,289.91)}, rotate = 180] [color={rgb, 255:red, 0; green, 0; blue, 0 }  ,draw opacity=1 ][fill={rgb, 255:red, 0; green, 0; blue, 0 }  ,fill opacity=1 ][line width=0.75]      (0, 0) circle [x radius= 3.35, y radius= 3.35]   ;
\draw [shift={(414,289.91)}, rotate = 180] [fill={rgb, 255:red, 0; green, 0; blue, 0 }  ,fill opacity=1 ][line width=0.08]  [draw opacity=0] (12,-3) -- (0,0) -- (12,3) -- cycle    ;
%Straight Lines [id:da8248821882706661] 
\draw [color={rgb, 255:red, 0; green, 0; blue, 0 }  ,draw opacity=1 ][fill={rgb, 255:red, 74; green, 144; blue, 226 }  ,fill opacity=1 ]   (424.56,251) -- (424.27,290) ;
\draw [shift={(424.27,290)}, rotate = 90.44] [color={rgb, 255:red, 0; green, 0; blue, 0 }  ,draw opacity=1 ][fill={rgb, 255:red, 0; green, 0; blue, 0 }  ,fill opacity=1 ][line width=0.75]      (0, 0) circle [x radius= 3.35, y radius= 3.35]   ;
\draw [shift={(424.58,249)}, rotate = 90.44] [fill={rgb, 255:red, 0; green, 0; blue, 0 }  ,fill opacity=1 ][line width=0.08]  [draw opacity=0] (12,-3) -- (0,0) -- (12,3) -- cycle    ;

% Text Node
\draw (94,76.4) node [anchor=north west][inner sep=0.75pt]  [font=\scriptsize]  {$( d+1,d-1)$};
% Text Node
\draw (185,100.4) node [anchor=north west][inner sep=0.75pt]  [font=\tiny]  {$d^{(2)}$};
% Text Node
\draw (162,102.4) node [anchor=north west][inner sep=0.75pt]  [font=\tiny]  {$\tilde{v}$};
% Text Node
\draw (227,116.4) node [anchor=north west][inner sep=0.75pt]  [font=\tiny]  {$\dbar$};
% Text Node
\draw (239,146.4) node [anchor=north west][inner sep=0.75pt]  [font=\tiny]  {$d^{(2)}$};
% Text Node
\draw (281,162.4) node [anchor=north west][inner sep=0.75pt]  [font=\tiny]  {$\dbar$};
% Text Node
\draw (284,201.4) node [anchor=north west][inner sep=0.75pt]  [font=\tiny]  {$d^{(2)}$};
% Text Node
\draw (326,217.4) node [anchor=north west][inner sep=0.75pt]  [font=\tiny]  {$\dbar$};
% Text Node
\draw (338,247.4) node [anchor=north west][inner sep=0.75pt]  [font=\tiny]  {$d^{(2)}$};
% Text Node
\draw (380,263.4) node [anchor=north west][inner sep=0.75pt]  [font=\tiny]  {$\dbar$};
% Text Node
\draw (481,278) node [anchor=north west][inner sep=0.75pt]   [align=left] {0};
% Text Node
\draw (420,308) node [anchor=north west][inner sep=0.75pt]   [align=left] {1};
% Text Node
\draw (367,308) node [anchor=north west][inner sep=0.75pt]   [align=left] {2};
% Text Node
\draw (432,263.4) node [anchor=north west][inner sep=0.75pt]  [font=\tiny]  {$\dbar$};
% Text Node
\draw (419,232) node [anchor=north west][inner sep=0.75pt]   [align=left] {0};
\end{tikzpicture}
                \caption{Zig-zag argument}
        \label{fig:Zig-Zag}
    \end{figure}
\end{proof}
\begin{table}[!h]
        \centering
\begin{tabular}{|p{0.12\textwidth}|p{0.12\textwidth}|p{0.12\textwidth}|p{0.12\textwidth}|p{0.12\textwidth}|}
\hline 
  & $\displaystyle \mathrm{ch}_{(1,1,1,1)}(3)$ & $\displaystyle \mathrm{ch}_{(3,1)}(3)$ & $\displaystyle\mathrm{ch}_{(2,2)}(3) $ & $\displaystyle \mathrm{ch}_{(4)}(4)$ \\
\hline 
 $\displaystyle \mathbb{v}_{(1,1,1,1)}( 3)$ & $\displaystyle \neq 0$ & $\displaystyle 0$ & $\displaystyle  0$ & $\displaystyle 0$ \\
\hline 
 $\displaystyle \mathbb{v}_{(3,1)}( 3)$ & $\displaystyle *$ & $\displaystyle \neq 0$ & $\displaystyle 0$ & $\displaystyle 0$ \\
\hline 
 $\displaystyle \mathbb{v}_{(2,2)}( 3)$ & $\displaystyle *$ & $\displaystyle *$ & $\displaystyle \neq 0$ & $\displaystyle 0$ \\
\hline 
 $\displaystyle \mathbb{v}_{(4)}( 3)$ & $\displaystyle *$ & $\displaystyle *$ & $\displaystyle *$ & $\displaystyle \neq 0$ \\
 \hline
\end{tabular}
\caption{Pairing matrix, $d=3$}
        \end{table}
        \begin{table}[!h]
        \centering
\begin{tabular}{|p{0.12\textwidth}|p{0.15\textwidth}|p{0.12\textwidth}|p{0.12\textwidth}|p{0.10\textwidth}|p{0.10\textwidth}|p{0.10\textwidth}|}
\hline 
  & $\displaystyle \mathrm{ch}_{(1,1,1,1,1)}(4)$ & $\displaystyle \mathrm{ch}_{(3,1,1)}(4)$ & $\displaystyle\mathrm{ch}_{(2,2,1)}(4) $ & $\displaystyle \mathrm{ch}_{(4,1)}(4)$ &$\displaystyle \mathrm{ch}_{(3,2)}(4)$& $\displaystyle \mathrm{ch}_{(5)}(4)$ \\
\hline 
 $\displaystyle \mathbb{v}_{(1,1,1,1,1)}( 4)$ & $\displaystyle \neq 0$ & $\displaystyle 0$ & $\displaystyle  0$ & $\displaystyle 0$& $\displaystyle 0$& $\displaystyle 0$ \\
\hline 
 $\displaystyle \mathbb{v}_{(3,1,1)}( 4)$ & $\displaystyle *$ & $\displaystyle \neq 0$ & $\displaystyle 0$ & $\displaystyle 0$ & $\displaystyle 0$& $\displaystyle 0$ \\
\hline 
 $\displaystyle \mathbb{v}_{(2,2,1)}( 4)$ & $\displaystyle *$ & $\displaystyle *$ & $\displaystyle \neq 0$ & $\displaystyle 0$& $\displaystyle 0$& $\displaystyle 0$  \\
\hline 
 $\displaystyle \mathbb{v}_{(4,1)}( 4)$ & $\displaystyle *$ & $\displaystyle *$ & $\displaystyle *$ & $\displaystyle \neq 0$ & $\displaystyle 0$& $\displaystyle 0$ \\
 \hline
 $\displaystyle \mathbb{v}_{(3,2)}( 4)$ & $\displaystyle *$ & $\displaystyle *$ & $\displaystyle *$ & $\displaystyle *$ & $\displaystyle \neq 0$& $\displaystyle 0$ \\
 \hline
  $\displaystyle \mathbb{v}_{(5)}( 4)$ & $\displaystyle *$ & $\displaystyle *$ & $\displaystyle *$ & $\displaystyle *$ & $\displaystyle *$& $\displaystyle \neq 0$ \\
 \hline
\end{tabular}
\caption{Pairing matrix, $d=4$}
\end{table}

For $I \subset \{1,\ldots,d\}$ we introduce the following notation
\begin{align}
  \mathrm{det}(z^{I},\mathsf{Eu}_{I})
& \define \sum\mathrm{sign}(\sigma)({z}^{\sigma(\ii_1)}\sum_{\ii\in I}z^{\ii}\partial_{\ii})\wedge\cdots\wedge ({z}
^{\sigma(\ii_l)}\sum_{\ii\in I}z^{\ii}\partial_{\ii}) \\
& =l!\cdot({z}^{\ii_1}\mathsf{Eu}_I)\wedge\cdots\wedge ({z}^{\ii_l}\mathsf{Eu}_I).
\end{align} 
Here, $\mathsf{Eu}_I$ is the Euler vector field in the directions $I$:
\begin{equation}\label{}
  \mathsf{Eu}_I \define \sum_{\ii \in I} z^{\ii} \del_{\ii} .
\end{equation}

Suppose that $\lambda = (\lambda_1,\ldots, \lambda_k)$, $\lambda_1 \geq \cdots \geq \lambda_k$ is a partition of the
integer $d+1$.
Such a $\lambda$ defines a set partition $\Lambda = (\Lambda_1,\ldots, \Lambda_k)$ of the set $\{0,\ldots,d\}$, with
$\lambda_i = |\Lambda_i|$ by
declaring: 
\begin{align*}
  \Lambda_1 & = \{0,1,\ldots, \lambda_1-1\} \\
  \Lambda_2 & = \{\lambda_1,\ldots, \lambda_1+\lambda_2-1\} \\
  & \vdots \\
  \Lambda_k & = \{\lambda_{k-1}, \ldots, d\} .
\end{align*}
When we write $\lambda$ for a partition of $d+1$ we always denote $\Lambda$ this particular set partition of
$\{0,\ldots,d\}$.
\begin{definition}\label{dfn:partitionv}
  For a partition $\lambda$ of $d+1$ with 
    \[
    \lambda_1 \geq \cdots \geq \lambda_k,\quad \lambda_1 \geq 2,
    \]
    define
    \[
      \mathbb{v}^{(d-1)}_{\lambda}(d) \define P(d) \partial_{\1}\wedge \mathbb{w}_\lambda(d) 
    \]
    where
    \begin{align*}
      \mathbb{w}_\lambda(d) &\define \sum\mathrm{sign}(\sigma)(z^{\1}{z}^{\sigma(\ii_1)}\mathsf{Eu}_{\Lambda_1-\{0\}}\wedge\cdots\wedge ({z}^{\sigma(\ii_l)}\mathsf{Eu}_{\Lambda_1-\{0\}})
                         \\ & \wedge \mathrm{det}(z^{\Lambda_2},\mathsf{Eu}_{\Lambda_2})\wedge\cdots \wedge\mathrm{det}(z^{\Lambda_{k-1}},\mathsf{Eu}_{\Lambda_{k-1}})\wedge\mathrm{det}(z^{\Lambda_{k}},\mathsf{Eu}_{\Lambda_{k}}).
    \end{align*}
    If $\lambda=(1,\dots,1)$, we define
  \begin{align*}
    \mathbb{v}^{(d-1)}_{(1,\dots,1)}(d) & \define P(d)\partial_{\1}\wedge \mathbb{w}_{(1,\dots,1)} (d) \\ & \; = P(d)
    \del_{\1} \wedge (z^{\1})^{3}\partial_{\1}\wedge (z^{\2})^{2}\partial_{\2}\wedge\cdots\wedge (z^{\dd})^{2}
    \partial_{\dd} .
  \end{align*}
\end{definition}

\begin{lemma}\label{lem:wlambda}
  For $d\geq2$ and $\lambda$ a partition of $\{1,\ldots,d\}$ one has:
    \begin{itemize}
      \item[(a)] The $\lie{w}_d$-chain $\mathbb{w}_{\lambda}(d)$ is $\bd^{(2)}$-closed.
      \item[(b)] There exists a $\lie{witt}_d$-chain $\mathbb{u}$ of total degree one such that $\bd^{(2)} \mathbb{v}
        ^{(d-1)}_{\lambda}(d) = \dbar
      \mathbb{u}$.
    \end{itemize}
\end{lemma}
\begin{proof}
  For part (a) we use the following elementary bracket identities
  \[
    [z^{\jj}\mathsf{Eu}_I,z^{\kkk}\mathsf{Eu}_I]=0 ,
  \]
  where $\jj,\kkk \in I$, and
  \[
    [z^\jj\mathsf{Eu}_I,z^\kkk\mathsf{Eu}_J]=0
  \]
  where $I\cap J=\varnothing$, $\jj\in I,\ \kkk\in J.$
  Every wedge factor appearing in $\mathbb{w}_\lambda(d)$ is of the form $z^{\jj}\mathsf{Eu}_I$ for one of the blocks $I=\Lambda_a$. 
  Thus any two wedge factors either belong to the same block, in which case the first identity applies, or to different blocks, in which case the second identity applies.

  Part (b). 
  Since $\mathbb{w}_\lambda (d)$ is $\bd^{(2)}$-closed, the only contributions to $\bd^{(2)}\mathbb{v}_\lambda^{(d-1)} (d)
  $ come from brackets of $P(d) \del_{\1}$ with the individual wedge factors of $\mathbb{w}_\lambda (d)$.
  Referring to the definition of $\op{det}(z^I, \mathsf{Eu}_I)$, it suffices to show the following type
  of bracket in $\lie{witt}_d$:
  \begin{equation}\label{}
    [P(d) \del_{\1} , f \mathsf{Eu}_I] 
  \end{equation}
  is $\dbar$ exact, where $f$ is an arbitrary homogeneous polynomial in $z$ which is at least linear.
  
  Using
  \[
    \mathsf{Eu}_I(P)=-d\,\tau_I P,
    \qquad
    \tau_I\define \sum_{\ii\in I}z^\ii x^{\ii},
  \]
  we compute
  \begin{equation}\label{eq:P-Eu-bracket-expanded}
    [P\del_\1,f\mathsf{Eu}_I]
    =
    P(\del_\1 f)\mathsf{Eu}_I
    +
    \mathbf 1_{\1\in I} fP\del_\1
    +
    d f\tau_I P\del_\1 .
  \end{equation}
  We argue that each term on the right-hand side is $\dbar$-exact. 
  The first and second terms are exact because their coefficients are holomorphic (depend only on $z$) positive-degree
multiples of $P$. The second term is exact for the same reason. 
For the third term, if $f=f(z)$ has degree $m>0$, then $f\tau_I P$ has Euler weight
$m-d$. 
Since $m \geq 1$, we see that $f\tau_I P$ is also $\dbar$-exact.
\end{proof}

By lemmas \ref{lem:descent} and \ref{lem:wlambda}, each class $\mathbb{v}^{(d-1)}_\lambda(d)$ lifts to a genuine cycle
\[
  \mathbb{v}_\lambda(d)
  \in C^{Lie}_\bullet(\lie{witt}_d)
\]
whose leading term is $\mathbb{v}^{(d-1)}_\lambda(d)$. 

\paragraph
Let
\[
  \op{Part}(\{0,\ldots,d\})^{-}
\]
denote the set of partitions of $\{0,\ldots,d\}$, ordered by block sizes, with
the partition of type $(2,1,\ldots,1)$ omitted. We order this set as follows.
First, we order partitions by increasing number of nonsingleton blocks, and
within each such stratum we use lexicographic order on the nonincreasing
sequence of block sizes. Equivalently, the order is chosen so that the examples
for $d=3$ and $d=4$ begin
\[
  (1,1,1,1)<(3,1)<(2,2)<(4)
\]
and
\[
  (1,1,1,1,1)<(3,1,1)<(2,2,1)<(4,1)<(3,2)<(5).
\]

\begin{proposition}\label{prop:upper-triangular-pairing}
  The pairing matrix
  \[
    \left(
    \big\langle \op{ch}_{\mu}(d),
    \til{\mathbb{v}}_{\lambda}(d)
    \big\rangle
    \right)_{\lambda,\mu\in \op{Part}(\{0,\ldots,d\})^{-}}
  \]
  is upper triangular with nonzero diagonal entries. More explicitly,
  \[
    \big\langle \op{ch}_{\lambda}(d),
    \til{\mathbb{v}}_{\lambda}(d)
    \big\rangle\neq 0
  \]
  and, if $\lambda<\nu$, then
  \[
    \big\langle \op{ch}_{\lambda}(d),
    \til{\mathbb{v}}_{\nu}(d)
    \big\rangle=0 .
  \]
\end{proposition}

\begin{proof}
  By lemmas \ref{lem:descent} and \ref{lem:wlambda}, we have shown that for each partition $\nu$ there is a cycle of the form
  \[
    \til{\mathbb{v}}_\nu(d)
    =
    \mathbb{v}^{(d-1)}_\nu(d)+\text{lower Jouanolou-degree terms}
  \]
  Importantly, the lower Jouanolou-degree terms do not contribute to the pairing with the Chern--Weil cocycles
  $\op{ch}_\lambda(d)$, for degree reasons. Hence it is enough to compute the
  pairing on the leading terms $\mathbb{v}^{(d-1)}_\nu(d)$.

  The leading term is
  \[
    \mathbb{v}^{(d-1)}_\nu(d)
    =
    P(d)\del_\1\wedge \mathbb{w}_\nu(d).
  \]
  Pairing the cocycle
  \[
    \op{ch}_\lambda
    =
    \op{ch}_1^{i_1}\cdots \op{ch}_d^{i_d}
  \]
  with this leading term amounts to distributing the trace factors among the wedge
  factors of $\mathbb{w}_\nu(d)$. 
  Each trace factor of length $r$ can be
  nonzero only if it is supported on a block of $\nu$ of size at least $r$.
  Moreover, the determinant expression
  \[
    \det(z^I,\mathsf{Eu}_I)
  \]
  forces the length-$|I|$ trace to use precisely the variables in the block
  $I$; if fewer or different variables are used, the alternating determinant
  gives zero.

  It follows that the pairing can be nonzero only when the block sizes required
  by $\lambda$ can be matched with the block sizes occurring in $\nu$. With the
  chosen order, this implies
  \[
    \lambda\leq \nu .
  \]
  Therefore
  \[
    \big\langle \op{ch}_{\lambda}(d),
    \mathbb{v}^{(d-1)}_{\nu}(d)
    \big\rangle=0
    \qquad \text{if } \lambda<\nu .
  \]

  On the diagonal, the matching is unique up to permutation of equal-sized
  blocks. In that case each determinant block is paired with the corresponding
  trace factor of the same length. The determinant antisymmetrization cancels
  the antisymmetrization in the trace and leaves a nonzero scalar multiple of
  \[
    \op{Res}\big(P(d)\,\d z^\1\cdots \d z^\dd\big) = 1
  \]
  Hence
  \[
    \big\langle \op{ch}_{\lambda}(d),
    \mathbb{v}^{(d-1)}_{\lambda}(d)
    \big\rangle\neq 0 .
  \]
  Since the correction terms do not affect the pairing, the same upper
  triangularity and nonvanishing diagonal hold for the cycles
  $\til{\mathbb{v}}_\lambda(d)$.
\end{proof}

\begin{theorem}\label{thm:cwd-injective}
  For every $d\geq 1$, the universal Chern--Weil homomorphism
  \[
    \mathsf{cw}_d\colon
    \C[\op{ch}_1,\ldots,\op{ch}_d]_{d+1}
    \to
    \HH^2_{\op{Lie}}(\lie{witt}_d)
  \]
  is injective.
\end{theorem}

\begin{proof}
  The case $d=1$ is classical, so assume $d\geq 2$.

  By Lemma \ref{lem:chd+1}, the class $\op{ch}_{d+1}$ lies in the span of the
  Chern--Weil classes
  \[
    \op{ch}_1^{i_1}\cdots \op{ch}_d^{i_d},
    \qquad
    \sum_{p=1}^d p\,i_p=d+1.
  \]
  Equivalently, we may replace the basis element corresponding to the partition
  $(2,1,\ldots,1)$ by the class $\op{ch}_{d+1}=\op{ch}_{(d+1)}$. Thus it is
  enough to prove linear independence of the classes
  \[
    \{\op{ch}_{\lambda}(d)\}_{\lambda\in
    \op{Part}(\{0,\ldots,d\})^{-}} .
  \]

  Suppose that there is a linear relation
  \[
    \sum_{\lambda\in \op{Part}(\{0,\ldots,d\})^{-}}
    a_\lambda\,\op{ch}_{\lambda}(d)=0
    \qquad
    \text{in }
    \HH^2_{\op{Lie}}(\lie{witt}_d).
  \]
  Pair this relation with the cycles
  $\til{\mathbb{v}}_\nu(d)$, ordered as in proposition
  \ref{prop:upper-triangular-pairing}. Let $\lambda_0$ be the smallest
  partition for which $a_{\lambda_0}\neq 0$. Pairing with
  $\til{\mathbb{v}}_{\lambda_0}(d)$ gives
  \[
    0
    =
    \sum_\lambda a_\lambda
    \big\langle
      \op{ch}_\lambda(d),
      \til{\mathbb{v}}_{\lambda_0}(d)
    \big\rangle .
  \]
  By upper triangularity, all terms with $\lambda>\lambda_0$ vanish, and by
  minimality all terms with $\lambda<\lambda_0$ have coefficient zero. Hence
  \[
    0
    =
    a_{\lambda_0}
    \big\langle
      \op{ch}_{\lambda_0}(d),
      \til{\mathbb{v}}_{\lambda_0}(d)
    \big\rangle .
  \]
  The diagonal pairing is nonzero, so $a_{\lambda_0}=0$, a contradiction.
  Therefore all $a_\lambda$ vanish, and the Chern--Weil classes are linearly
  independent.

  Hence $\mathsf{cw}_d$ is injective.
\end{proof}

\subsection{Injectivity for the $\lie{gl}_r$-extended Chern--Weil homomorphism}

Fix $0 \leq k \leq d$.  Let
\[
  \nu=(1^{i_1},2^{i_2},\ldots,(d-k)^{i_{d-k}})
\]
be a partition of $d-k$, and let $\kappa=(k_1,\ldots,k_r)$ satisfy
\[
  k_1+2k_2+\cdots+rk_r=k+1 .
\]
We write
\begin{equation}\label{eq:ch-nu-kappa}
  \op{ch}_{\nu;\kappa}(d,k)
  \define
  \op{ch}_1^{i_1}\cdots \op{ch}_{d-k}^{i_{d-k}}
  \lie{t}_1^{k_1}\cdots \lie{t}_r^{k_r}.
\end{equation}

Let
\[
  \theta_\kappa(A)
  \define
  \op{Tr}(A)^{k_1}\op{Tr}(A^2)^{k_2}\cdots\op{Tr}(A^r)^{k_r}.
\]
The polynomials $\theta_\kappa$, with $\sum_j j k_j=k+1$, form a basis of
$(S^{k+1}\lie{gl}_r^*)^{\lie{gl}_r}$.  Choose the dual basis
\[
  \mathbb{E}_\kappa \in (S^{k+1}\lie{gl}_r)^{\lie{gl}_r},
  \qquad
  \langle \theta_\kappa,\mathbb{E}_{\kappa'}\rangle=\delta_{\kappa,\kappa'}.
\]
Write
\[
  \mathbb{E}_\kappa
  =
  \sum_{\ell}
  A^{\kappa,\ell}_0\odot A^{\kappa,\ell}_1\odot\cdots\odot A^{\kappa,\ell}_k .
\]

\begin{definition}\label{dfn:partitionvv}
Let $\lambda=(\lambda_1,\ldots,\lambda_{d-k})$ be a partition of $d-k$, written with
$\lambda_1\geq\cdots\geq\lambda_{d-k}$.  Define the chain
\[
\begin{aligned}
\mathbb{v}_{\lambda,\kappa}(d)
\define
\sum_{\tau\in S_{k+1}}\sum_{\ell}
&
P(d)\,
A^{\kappa,\ell}_{\tau(0)}
\wedge z^{\1}A^{\kappa,\ell}_{\tau(1)}
\wedge\cdots\wedge
z^{\kkk}A^{\kappa,\ell}_{\tau(k)}
\\
&\wedge
\det(z^{\Lambda_1},\mathsf{Eu}_{\Lambda_1})
\wedge\cdots\wedge
\det(z^{\Lambda_{d-k}},\mathsf{Eu}_{\Lambda_{d-k}}).
\end{aligned}
\]
\end{definition}
\begin{lemma}\label{lem:wlambdaMixed}
For $d\geq 2$, the chain $\mathbb{v}_{\lambda,\kappa}(d)$ satisfies
\[
  \bd^{(2)}\mathbb{v}_{\lambda,\kappa}(d)
  =
  \dbar \mathbb{u}
\]
for some $\lie{witt}_d\ltimes \lie{gl}_r(\cJ_d)$-chain $\mathbb{u}$ of total degree one.
\end{lemma}

\begin{proof}
The $\lie{gl}_r(\cJ_d)$-part
\[
  \sum_{\tau\in S_{k+1}}\sum_{\ell}
  P(d)\,
  A^{\kappa,\ell}_{\tau(0)}
  \wedge z^{\1}A^{\kappa,\ell}_{\tau(1)}
  \wedge\cdots\wedge z^{\kkk}A^{\kappa,\ell}_{\tau(k)}
\]
is $\bd^{(2)}$-closed, because $\cJ_d$ is commutative.  The remaining factors are the same Witt chains as in
lemma~\ref{lem:wlambda}, so the same descent argument gives the desired $\mathbb{u}$.
\end{proof}

\begin{theorem}\label{thm:cwdr}
The $\lie{gl}_r$-extended universal Chern--Weil map
\[
  \mathsf{cw}_{d,r}
  \colon
  \C[y_1,\ldots,y_d,\lie{t}_1,\ldots,\lie{t}_r]_{d+1}
  \longrightarrow
  \HH^2_{\op{Lie}}
  \bigl(\lie{witt}_d\ltimes \lie{gl}_r(\cJ_d)\bigr)
\]
is injective.
\end{theorem}

\begin{proof}
Pair the classes $\op{ch}_{\nu;\kappa}(d,k)$ with the chains
$\mathbb{v}_{\lambda,\kappa'}(d)$.  By Lemma~\ref{lem:wlambdaMixed}, these chains can be lifted to cycles, so the pairings are homology invariants.

The $\lie{gl}_r$-factor separates the different $\kappa$'s:
\[
  \big\langle \theta_\kappa,\mathbb{E}_{\kappa'}\big\rangle
  =
  \delta_{\kappa,\kappa'}.
\]
Thus the pairing matrix is block diagonal with respect to $\kappa$.  Inside each block, the Witt part is exactly the
pairing matrix appearing in the proof of lemma~\ref{lem:wlambda}; with the same ordering on partitions, it is upper triangular with nonzero diagonal entries.  Therefore the full pairing matrix is upper triangular with nonzero diagonal.  Hence no nontrivial linear combination of the classes
$\op{ch}_{\nu;\kappa}(d,k)$ lies in the kernel of $\mathsf{cw}_{d,r}$, proving injectivity.
\end{proof}
\end{document}

